\section{Detailed Proof Sketches}

This appendix collects the longer derivation steps omitted from the main text. The goal is not to expand
the model beyond the submission claim, but to show how the topology-level formulas in Section~5 follow
from the visibility assumptions.

\subsection{Reviewer Gates and Correlated Failure}

Let $X = \mathbf{1}_{E_{\mathrm{gen}}}$ and $Y = \mathbf{1}_{M_{\mathrm{ver}}}$. In the direct topology,
failure is exactly $E_{\mathrm{gen}}$, so
\[
P(\mathrm{Fail}_{\mathrm{direct}})=P(E_{\mathrm{gen}})=p.
\]
In the reviewer-gated topology, failure is the conjunction
$E_{\mathrm{gen}} \wedge M_{\mathrm{ver}}$. Under independence this gives
\[
P(\mathrm{Fail}_{\mathrm{gate}})=pq.
\]
For the correlated case,
\[
P(E_{\mathrm{gen}} \wedge M_{\mathrm{ver}})
= \mathbb{E}[XY]
= \mathbb{E}[X]\mathbb{E}[Y] + \mathrm{Cov}(X,Y)
= pq + \rho\sqrt{p(1-p)q(1-q)}.
\]
The topology contributes the conjunction; implementation choices determine the covariance term.

\subsection{Error Amplification}

Let $E_j$ denote the event that worker $j$ emits an erroneous output, with $P(E_j)=p$ independently
across workers.

\paragraph{Independent aggregation.}
The aggregate fails when at least one worker error survives to the combined output. The exact failure
probability is
\[
P(\mathrm{failure}_{\otimes}) = 1 - (1-p)^n.
\]
Applying the union bound gives
\[
P(\mathrm{failure}_{\otimes}) \le \sum_{j=1}^{n} P(E_j) = np,
\]
so the amplification factor satisfies $A_{\otimes} \le n$.

\paragraph{Centralized aggregation.}
Suppose a hub observes all worker outputs before release and detects a worker error with rate
$d = P(\text{hub detects error}\mid \text{error occurred})$. A worker error reaches the final output only
if it occurs and escapes detection, so the exact centralized failure probability is
\[
P(\mathrm{failure}_{\circ}) = 1 - (1-p(1-d))^n.
\]
Again applying the union bound yields
\[
P(\mathrm{failure}_{\circ}) \le np(1-d),
\]
hence $A_{\circ} \le n(1-d)$.

\paragraph{Exact comparison.}
The ratio
\[
\frac{A_{\otimes}}{A_{\circ}}
=
\frac{1-(1-p)^n}{1-(1-p(1-d))^n}
\xrightarrow[p \to 0]{} \frac{1}{1-d}.
\]
shows that the small-$p$ regime recovers the intuitive $1/(1-d)$ gain.

\subsection{Sequential Handoff Overhead}

Let step $j$ be executed by agent $a$ and step $j+1$ by agent $b$. The sender's full task-relevant state
is $S_a$, while the receiver sees only a transmitted summary $\mathrm{obs}_b(S_a)$.

Applying the data processing inequality to the Markov chain
$S_a \to \mathrm{obs}_b(S_a) \to \hat r_j$ gives
\[
I(\mathrm{obs}_b(S_a); r_j) \le I(S_a; r_j).
\]
Define the lost task-relevant information as
\[
\Delta I_j = I(S_a; r_j) - I(\mathrm{obs}_b(S_a); r_j) \ge 0.
\]
Equality holds only if the handoff is lossless for the task-relevant signal. In practical agent systems,
finite context windows, summarization, and schema compression make strict equality atypical.

If $h$ handoffs occur in a $k$-step task, the total multi-agent cost can be written as
\[
C_{\mathrm{multi}}
=
k \cdot c_{\mathrm{step}} + \sum_{j=1}^{h} \bigl(c_{\mathrm{comm}} + c_{\mathrm{recon},j}\bigr),
\]
where $c_{\mathrm{recon},j}$ is the receiver's effort to compensate for the missing context. This yields
the overhead ratio used in the main text and explains why repeated handoffs on strictly sequential tasks
can accumulate superlinearly in practice.

\subsection{Parallel Acceleration}

Suppose a task decomposes into subtasks $\varphi_1,\dots,\varphi_m$ and that each worker can solve its
assigned subtask from local observations alone. Let $c_{\mathrm{sub}_i}$ be the wall-clock cost of
subtask $i$.

\paragraph{Parallel cost.}
When all subtasks run concurrently, total wall-clock time is the slowest subtask plus coordinator
assignment and aggregation overhead:
\[
C_{\mathrm{parallel}} = \max_i(c_{\mathrm{sub}_i}) + c_{\mathrm{assign}} + c_{\mathrm{agg}}.
\]
The single-agent sequential baseline is
\[
C_{\mathrm{sequential}} = \sum_{i=1}^{m} c_{\mathrm{sub}_i},
\]
which yields the speedup formula used in the main text.

\paragraph{Coordinator visibility.}
If the coordinator must receive the tuple of worker outputs $(o_1,\dots,o_m)$ to aggregate them, then
for any proposition $\varphi$ measurable with respect to that tuple, pooled worker-output knowledge
implies coordinator knowledge:
\[
D_G(\varphi) \Rightarrow K_{\mathrm{hub}}(\varphi).
\]
This is why cross-checking comes ``for free'' in centralized decompositions: the aggregation step already
creates pooled visibility.

\paragraph{Epistemic ceiling.}
If some worker's optimal action depends on another worker's hidden local result, then the subtasks are not
epistemically independent. Running them in parallel forces a worker to act without a needed fact, which
appears in the main text as quality loss proportional to the unresolved mutual information.

\subsection{Tool Density}

Let $T_1,\dots,T_n$ be an approximately balanced partition of $t$ tools across $n$ agents, so
$|T_i| \approx t/n$.

\paragraph{Remote-tool coordination.}
If agent $i$ needs tool $k \in T_j$ with $j \ne i$, three steps are required:
\begin{enumerate}[leftmargin=*]
\item establish that some remote tool can solve the current subproblem;
\item delegate execution to the remote owner;
\item reconstruct and interpret the returned result without the owner's full local context.
\end{enumerate}
This is the operational content of the coordination term in
the main text.

\paragraph{Second-order planning.}
Planning is no longer just ``which tool should I call?'' but
``which agent knows how to apply which tool to the current subproblem?'' In epistemic terms, the planner
must reason about propositions of the form
$K_i(K_j(\mathrm{tool}_k \text{ can solve } \varphi))$.
In the worst case this introduces $O(tn)$ comparisons across tool-agent pairs, versus $O(t)$ for the
single-agent baseline.
